\subsection{FINITE NILPOTENT GROUPS}

THEOREM 4. Let G be a finite group. The following conditions are equivalent:

\begin{enumerate}
\def\labelenumi{(\alph{enumi})}
\item
  G is nilpotent.
\item
  G is a product of p-groups.
\item
  For every prime number \(p\) there exists a normal Sylow \(p\) -subgroup of \(G\) .
\item
  (b) \(\Rightarrow\) (a) (no. 5, Corollary to Theorem 1).
\end{enumerate}

Suppose (a) holds and let P be a Sylow p-subgroup of G. If N is the normalizer of P in G, Corollary 1 to Theorem 3 shows that N is its own normalizer. By §6, no. 3, Corollary to Proposition 8, this shows that N = G. Hence (a) ⇒ (c).

Suppose (c) holds and let I be the set of prime numbers dividing \(\operatorname{Card}(G)\) . For all \(p \in I\) , let \(P_p\) be a normal Sylow \(p\) -subgroup of \(G\) . For all \(p \neq q\) , \(P_p \cap P_q\) is reduced to \(e\) for it is both a \(p\) -group and a \(q\) -group, hence \(P_p\) and \(P_q\) centralize one another (§ 4, no. 9, Proposition 15). Let \(\phi\) be the canonical homomorphism (§ 4, no. 9, Proposition 12) of \(\prod_{p \in I} P_p\) into \(G\) . The homomorphism \(\phi\) is surjective by the Remark of no. 6. As \(\operatorname{Card}\left(\prod_{p \in I} P_p\right) = \operatorname{Card}(G)\) , it follows that \(\phi\) is bijective.

Remarks. (1) Let G be a finite group and p a prime number. By no. 6, Theorem 3 (a) and no. 6, Theorem 2, the following conditions are equivalent:

\begin{enumerate}
\def\labelenumi{(\roman{enumi})}
\item
  there exists a normal Sylow \(p\) -subgroup of \(\mathbf{G}\) ;
\item
  every Sylow \(p\) -subgroup of \(\mathbf{G}\) is normal;
\item
  there exists only one Sylow p-subgroup of G.
\end{enumerate}

\begin{enumerate}
\def\labelenumi{(\arabic{enumi})}
\setcounter{enumi}{1}
\item
  Let \(\mathbf{G}\) be a nilpotent finite group. Let \(\mathbf{I}\) be the set of prime divisors of \(\operatorname{Card}(\mathbf{G})\) . By Theorem 4 and Remark 1, \(\mathbf{G} = \prod_{p \in \mathbf{I}} \mathbf{G}_p\) , where \(\mathbf{G}_p\) is the unique Sylow \(p\) -group of \(\mathbf{G}\) .
\item
  Applied to commutative groups, Theorem 4 gives the decomposition, of commutative finite groups as a product of primary components, which will be studied from another point of view in Chapter VII.
\end{enumerate}

Example. The group \(\mathfrak{S}_3\) is of order 6. It contains a normal Sylow 3-subgroup of order 3: the group \(\mathfrak{A}_3\) . It contains three Sylow 2-subgroups of order 2: the groups \(\{e, \tau\}\) , where \(\tau\) is a transposition. The group \(\mathfrak{S}_3\) is thus not nilpotent.

